%==============================================================================
%  09-conclusion.tex — Section 8: Executive Conclusion & Call to Action
%==============================================================================
\strategyPart{24 \textbar\ Executive Conclusion}{Executive Conclusion and Call to Action}

The evidence assembled in this document converges on a single strategic
imperative: \textbf{the organization that treats AI as a problem-solving
capability will compound value; the organization that treats AI as a tool
portfolio will accumulate cost.}

\section{The strategic imperative in numbers}

\vspace{-0.3em}
\begin{center}
  \stat{\$13T}{AI value by 2030\\(McKinsey 2018)}\hfill
  \stat{\$4.4T}{Annual gen-AI value\\(McKinsey 2023)}\hfill
  \stat{65\%}{Of orgs now use gen AI\\(McKinsey 2024)}
\end{center}

\begin{center}
  \stat{85\%}{Of AI projects errored\\(Gartner 2018)}\hfill
  \stat{30\%}{Abandoned after PoC\\(Gartner 2023)}\hfill
  \stat{7/10}{Report minimal impact\\(MIT SMR--BCG 2019)}
\end{center}

The first row is the size of the prize; the second row is the probability of
mismanaging it. Both are governed by the same variable: whether the
organization leads with problems or with technology.

\section{Three immediate executive asks}

\begin{enumerate}
  \item \textbf{Authorize the audit of the organizational ``stocks.''} Direct
        a ninety-day assessment of human capital, data capital and
        organizational capital against the maturity stages in Section 4. The
        audit's output is a readiness scorecard, not a technology roadmap
        [MIT CISR].
  \item \textbf{Convene an AI Reinvention Council.} Establish a cross-
        functional coalition — strategy, operations, data, HR and finance —
        with the mandate to own the problem-first agenda end to end, mirroring
        the guiding-coalition step of Kotter's model [Kotter Inc.].
  \item \textbf{Commit to a problem-first strategy.} Require that every AI
        investment be filed against a named organizational problem, a measured
        baseline, a KPI and an accountable owner — and that no model be funded
        on the strength of a technology alone [McKinsey 2023; Gartner 2023].
\end{enumerate}

\vspace{1em}
\begin{center}
  \yBanner[color=inBright]{We do not implement AI for the organization; we solve organizational problems using AI.}
\end{center}

\vspace{0.5em}
\noindent The window for action is open, but it is not open indefinitely. The
economic prize is measured in trillions, the failure rate is measured in the
majority, and the frontier — agentic AI — is arriving in years, not decades
[Gartner 2024]. The board and the executive committee should treat this
document not as a briefing but as a decision: to close the organizational gap
now, or to compete from behind.

\section{The design principle for AI agents}

The strongest design principle that synthesizes this document:

\begin{center}
  \yBanner[color=inNavy]{Use assistants for judgment-heavy work, workflows for predictable processes,\\agents for bounded and measurable tasks, and multi-agent systems only when\\the problem genuinely requires dynamic coordination among specialists.}
\end{center}

Organizations should not replace people with agents based only on technical
capability. They should redesign work by separating tasks into activities AI
can automate, activities AI can assist with, and responsibilities that must
remain human-led because they require judgment, accountability, empathy,
ethical reasoning or legal responsibility [Anthropic; NIST AI RMF; Microsoft
agent framework].

\takeaway{Audit the three stocks, convene an AI Reinvention Council, and commit to a problem-first strategy.}
